
\usepackage{mathpartir}

\usepackage{graphicx}

% \usepackage[firstpage]{draftwatermark}
% % \SetWatermarkText{\hspace*{8in}\raisebox{3.5in}{\includegraphics[scale=0.1]{aec-badge-popl.pdf}}}
% \SetWatermarkText{\hspace*{13.3in}\raisebox{-6.5in}{\includegraphics[scale=0.1]{aec-badge-popl.pdf}}}
% \SetWatermarkAngle{0}

%--------------------------------------
% To use \mathbb
\usepackage{amsfonts}
%--------------------------------------
% To defined theorems and Lemmas
 \usepackage{amsthm} 
\theoremstyle{definition}  % non-italic
% Definition, Lemma and Theorem
\newtheorem{definition}{Definition}
\newtheorem{theorem}{Theorem}
% \newtheorem{definition}{Lemma}
\newtheorem{lemma}{Lemma}
\newtheorem{corollary}{Corollary}
%--------------------------------------
% To have colored annotations
\usepackage{xcolor}
%--------------------------------------
% To strike out text
\usepackage[normalem]{ulem}
\usepackage{cancel}
%--------------------------------------
%\usepackage{hyperref}
\usepackage{url}
%--------------------------------------
% Programming language rules
\usepackage{mmm}
\usepackage{pl}
%--------------------------------------
% Arrow symbols 
\let\Bbbk\relax
\usepackage{amssymb}
%--------------------------------------
% \xrightarrow
\usepackage{amsmath}
%--------------------------------------
% % To write algorithms
% \usepackage{algorithm}
% %\usepackage{algpseudocode}
% \usepackage[noend]{algpseudocode}
% \let\algState\State
% \usepackage{varwidth}% http://ctan.org/pkg/varwidth
% % To have Algorithm float
% % \usepackage{algorithm}
% % To make the algorithms float to the top
% \usepackage{float}
% \newfloat{algorithm}{t}{lop}
% %\renewcommand{\algorithmicforall}{\textbf{foreach}}
%--------------------------------------

% this employs smarter hyphenations, par, and other rules
% ... and helps condense the paper
\usepackage{microtype}

% To remove a warning for fonts
% \usepackage{lmodern}% http://ctan.org/pkg/lm

%--------------------------------------
% Symbol definitions and comments
% Standard
\def\t{\phantom{X}}

\def\<{\langle}
\def\>{\rangle}

%--------------------------------------

\usepackage{listings}
\lstset{
   language=C,
   captionpos=b,
   tabsize=3,
%   frame=lines,
%   keywordstyle=\color{blue},
%   commentstyle=\color{darkgreen},
%   stringstyle=\color{red},
   numbers=left,
   numberstyle=\tiny,
   numbersep=5pt,
   breaklines=true,
   showstringspaces=false,
%   basicstyle=\small\bfseries,
   basicstyle=\small\ttfamily,
   emph={label}
}

\usepackage{array}

\usepackage{tikz}
\usetikzlibrary{positioning} 
\usetikzlibrary{arrows} 
%\usetikzlibrary{graphs} % Not supported by tikz 2.10
\usetikzlibrary{calc}
\usetikzlibrary{backgrounds}
\usetikzlibrary{shapes}

%\usepackage{caption}
%\usepackage{subcaption}

\usepackage{stmaryrd}

\def\defeq{\triangleq}

\def\self{\mathit{self}}

\def\switch{\mathsf{switch} \ }
\def\case{\mathsf{case} \ }
\def\ret{\mathsf{ret} \ }
\def\lforall{\mathsf{forall} \ }
\def\lif{\mathsf{if} \ }
\def\lelse{\mathsf{else} \ }
\def\lmap{\mathsf{map} \ }
\def\llet{\mathsf{let} \ }
\def\max{\mathsf{max}}

\def\vc{\mathit{vc}}
\def\List{\mathsf{List}}
\def\llookup{\mathsf{lookup} \ }

\def\RYW{\mathit{RYW}}
\def\MR{\mathit{MR}}
\def\WFR{\mathit{WFR}}
\def\MW{\mathit{MW}}
\def\CC{\mathit{CC}}
\def\LCC{\mathit{LCC}}
\def\Rel{\mathit{Rel}}

\def\ext{\mathsf{ext}}

\def\CState{\mathsf{CState}}
\def\RState{\mathsf{RState}}

\def\true{\mathsf{true}}
\def\Set{\mathsf{Set}}
\def\Nat{\mathsf{Nat}}
% \def\L{\mathsf{L}}
\def\L{L}
\def\lab{\mathit{label}}
\def\clientLabs{\mathit{client\text{-}labels}}

\def\cInit{\mathsf{c\text{-}init}}
\def\rInit{\mathsf{r\text{-}init}}

\def\get{\mathsf{get}}
\def\getGuard{\mathsf{get\text{-}guard}}
\def\getReq{\mathsf{get\text{-}req}}
\def\getRes{\mathsf{get\text{-}res}}

\def\lput{\mathsf{put}}
\def\putGuard{\mathsf{put\text{-}guard}}
\def\putReq{\mathsf{put\text{-}req}}
\def\putRes{\mathsf{put\text{-}res}}

\def\Type{\mathsf{Type}}
\def\Option{\mathsf{Option}}
\def\Bool{\mathsf{Bool}}

\def\Event{\mathit{Event}}
\def\Req{\mathit{Req}}
\def\req{\mathit{req}}

\def\Res{\mathit{Res}}
\def\res{\mathit{res}}

\def\PutReq{\mathit{Put\text{-}req}}
\def\PutReqPayload{\mathsf{Put\text{-}req\text{-}payload}}
\def\PutRes{\mathit{Put\text{-}res}}
\def\PutResPayload{\mathsf{Put\text{-}res\text{-}payload}}

\def\GetReq{\mathit{Get\text{-}req}}
\def\GetReqPayload{\mathsf{Get\text{-}req\text{-}payload}}
\def\GetRes{\mathit{Get\text{-}res}}
\def\GetResPayload{\mathsf{Get\text{-}res\text{-}payload}}

\def\Unit{\mathsf{Unit}}

% Program Syntax

% Uses:
% $
% \sPut{}{},
% \sPut{k}{v},
% \sPut[i]{k}{v},
% % 
% \sGet{}{},
% \sGet{x}{k},
% \sGet[i]{x}{k},
% $

\newcommand\sPut[3][]{%
  \mathit{put}%
  \ifthenelse{\equal{#1}{}}{%
    \ifthenelse{\equal{#2}{}}{}{
    ({#2},{#3})}%
  }{%
    ^{#1}({#2},{#3})%
  }%
}  
  
\newcommand\sGet[3][]{%
   \ifthenelse{\equal{#2}{}}{}{#2\leftarrow}%
   \mathit{get}
   \ifthenelse{\equal{#1}{}}{}{^{#1}}%
   \ifthenelse{\equal{#3}{}}{}{({#3})}}

% \newcommand\sGet[2][]{%
%     \ifthenelse{\equal{#1}{}}{}{#1\leftarrow }%
%     \mathit{get}%
%     \ifthenelse{\equal{#2}{}}{}{({#2})}}
    
  \newcommand\sSkip{\mathit{skip}}

\newcommand{\name}[1]{\textsc{#1}}

\newcommand\indrule[3][]{%
%  \inferrule*[Right=#1]{#2}{#3}%
  \inferrule[#1]{#2}{#3}%
  }

\newcommand\figurealgtextsize{%
%   \small%
  \setlength{\abovedisplayskip}{0pt}%
  \setlength{\abovedisplayshortskip}{0pt}%
  \setlength{\belowdisplayskip}{0pt}%
  \setlength{\belowdisplayshortskip}{0pt}}
\newcommand\figureruletextsize{%
%   \small
  \setlength{\abovedisplayskip}{0pt}%
  \setlength{\abovedisplayshortskip}{0pt}%
  \setlength{\belowdisplayskip}{0pt}%
  \setlength{\belowdisplayshortskip}{0pt}
}
\newcommand\figuredefstextsize{%
%   \small
%  \large  
  \setlength{\abovedisplayskip}{0pt}%
  \setlength{\abovedisplayshortskip}{0pt}%
  \setlength{\belowdisplayskip}{0pt}%
  \setlength{\belowdisplayshortskip}{0pt}
}

\usepackage{siunitx}  % To format numbers

\setlength{\textfloatsep}{20pt}

\usepackage{balance}

% \newfont{\mycrnotice}{ptmr8t at 7pt}
% \newfont{\myconfname}{ptmri8t at 7pt}
% \let\crnotice\mycrnotice%
% \let\confname\myconfname%

\newcommand\world[3]{%
   {\begin{array}{l} ( %\left(
      #1, \\ \phantom{(}
      #2, \\ \phantom{(}
      #3          
   ) \end{array}}
}  

\newcommand\note[1]{\text{\small#1}}
\newcommand\bnote[1]{\textbf{\small#1}}

% \newcommand\note[1]{\footnotesize\text{#1}}
% \newcommand\bnote[1]{\footnotesize\textbf{#1}}

\usepackage{tikz}
\usetikzlibrary{positioning, arrows.meta}
